LUT factorisation reduces the query-specific table; code layout determines
how lanes fetch each candidate's table indices. These are separate
optimisations. The scan assigns neighbouring lanes to neighbouring
candidate occurrences, so its broad phase benefits from storing codes
by subspace across candidates rather than storing each candidate's full
primary code contiguously.

\subsection{Address mapping and four-byte packing}

Let $o_\ell$ be the start of inverted list $\ell$, $r$ the row within
that list and $p=o_\ell+r$ its physical position. If $P[m,p]$ denotes
the eight-bit code for primary subspace $m$, the byte-oriented scan copy
stores it at byte offset $mN+p$. Candidate-major storage would instead
use $pM+m$. When lanes visit consecutive positions within a list, the
former makes their accesses at fixed $m$ adjacent. At list boundaries,
positions can jump, so the mapping does not guarantee one contiguous
transaction for every warp.

The packed representation combines four subspace codes of the same
candidate into one unsigned word:
\begin{equation}
 W[j,p]=\sum_{t=0}^{3}2^{8t}P[4j+t,p],\qquad
 \operatorname{offset}_{\rm bytes}(W[j,p])=4(jN+p).
 \label{eq:5-packed-layout}
\end{equation}
The four codes are recovered exactly by
$(W[j,p]\gg(8t))\mathbin{\&}255$. Each recovered code then supplies the
high and low nibble indices in \cref{eq:5-lut-factorisation}. Packing
four subspaces is distinct from Cartesian splitting of one subspace's
eight-bit label, and from packing codes for four different candidates.

\begin{figure}[tbp]
\centering
% Address schematic; four lanes stand for a subset of a warp.
\begin{tikzpicture}[font=\small,
 box/.style={draw=ThesisNavy!65,rounded corners=2pt,align=center,
   text width=2.6cm,minimum height=1.2cm,inner sep=4pt},
 arr/.style={-{Latex[length=2mm]},draw=ThesisNavy!80,thick}]
\foreach \a in {0,1,2,3} {
 \node[box,fill=ThesisNavy!5] (w\a) at (3.3*\a,0)
  {lane $\a$\\$\mathsf{w}[j,p+\a]$};
}
\node[align=center] at (4.95,1.05)
 {Fixed $j$: adjacent 32-bit words across candidate positions};
\node[box,text width=11.5cm,minimum height=1.6cm] (decode) at (4.95,-2.5)
 {Unpack each lane's word independently at its position $p_a$\\
  $Z[4j,p_a]\;|\;Z[4j+1,p_a]\;|\;Z[4j+2,p_a]\;|\;Z[4j+3,p_a]$\\
  Four code bytes $\longrightarrow$ shifts and masks $\longrightarrow$ LUT indices};
\draw[arr] (w0.south) -- (decode.north west);
\draw[arr] (w1.south) -- (decode.north);
\draw[arr] (w2.south) -- (decode.north);
\draw[arr] (w3.south) -- (decode.north east);
\end{tikzpicture}

\caption{Packed candidate-parallel access, showing a subset of warp lanes.
At fixed word group $j$, consecutive positions access adjacent four-byte
words. Each word contains four subspace codes of one candidate. Unpacking
preserves each code byte; LUT splitting happens afterwards. The diagram
illustrates addresses, not a measured number of memory transactions.}
\label{fig:5-packed-access}
\end{figure}

The build path explicitly repacks the $[M,N]$ byte array into
$[M/4,N]$ words. Merely reinterpreting the original byte array as
\texttt{uint32\_t} would combine consecutive candidate bytes rather
than the required subspaces. The implementation requires $M$ divisible
by four and rejects a partial final word. The repack reads and writes
$NM$ code bytes once at construction, with old and new storage briefly
coexisting; it is outside the search timing boundary.

\subsection{What packed access saves, and what it preserves}

At the source level, a candidate now obtains four codes with one word
load instead of four byte loads, followed by shifts and masks in
registers. Within a contiguous list segment, 32 lanes request a logical
128-byte word window. Actual transactions depend on alignment, active
lanes and the memory hierarchy. A fourfold reduction in source-level
load operations does not establish a fourfold reduction in transactions
or elapsed time.

The logical payload remains $M$ bytes per candidate: this repacking
does not further compress the index. Both byte and word alternatives
already organise the scan copy by subspace or subspace group. The
historical on/off comparison therefore measures the net effect of the
packed layout and its load/unpack path, not the isolated contribution
of transposing candidate-major storage. It also cannot uniquely separate
instruction count, addressing and cache effects.

\Cref{sec:7-3} reports 48 matched byte/word configurations with throughput
ratios $1.012$--$1.658\times$. Each arm has one aggregate rate, so these
records provide weaker variability evidence than the newer repeated
LUT sweep. The integer representation identity is exact, while the
observed recall spread of at most $10^{-4}$ belongs to the complete
measured execution. The method claim is preservation of code values
with fewer load/address operations, qualified by the experiment's
measurement scope.
